\section{Absorbing localized logarithms into the forced gas}
\label{sec:soft-log-fusion}
The two exponential shifts in \eqref{eq:leaf-tilted-identity} arise
from an exact regrouping of the forced activities with localized
logarithms. We prove this regrouping, including its flag constraints
and the derivative price needed for a tilted family.

All activities and logarithms in this section are real valued.
Fix a coarse configuration $U$ and let $F(U)$ be its set of flagged
sites. A hard activity $K_i(U)$ has a nonempty support $X_i$,
a nonempty inventory $I_i\subseteq X_i$, and a fixed eligibility
condition. The hard factors and their eligibility predicates are
measurable, gauge invariant and local on $X_i$, with the same
final-flag convention as the logarithms below. Nonzero hard values
satisfy their eligibility condition. The hard gas
sums over finite families $A$ of pairwise disjoint supports with
$\bigcup_{i\in A}I_i=F(U)$:
\[
                    G(U)=\sum_{A\ {\rm admissible}}\prod_{i\in A}K_i(U).
\]
The empty family is allowed precisely when $F(U)=\varnothing$.
Each logarithmic term $\ell_j(U)$ has nonempty support $Y_j$,
is measurable and gauge invariant, and is local with the same
threshold convention as the hard activities. Assume
\begin{equation}\label{eq:fusion-vanishing}
                  Y_j\cap F(U)=\varnothing\ \Longrightarrow\ \ell_j(U)=0.
\end{equation}
Suppose there are summable uniform majorants for both $K_i$ and
$\ell_j$. These hypotheses will also imply absolute convergence of
the new gas; it need not be assumed separately.

Use a disjoint union of the hard and logarithmic labels, assigning
the factors
\[
                        r_i=K_i,\qquad r_j=e^{\ell_j}-1
\]
and the supports $X_i$, $Y_j$, respectively. A new activity label
$C$ is a finite connected set of these labels which contains at
least one hard label and whose hard supports are pairwise disjoint.
Write $Z_C$ for the union of all its supports and
$J_C=\bigcup_{i\in C\ {\rm hard}}I_i$ for its inventory. The new
activity is
\begin{equation}\label{eq:fusion-new-activity}
 \widehat K_C(U)=
 \mathbf1_{\{Z_C\cap F(U)=J_C\}}
 \mathbf1_{\{\text{all hard labels in }C\text{ are eligible at }U\}}
                         \prod_{v\in C}r_v(U).
\end{equation}
Its support and inventory are fixed by $C$ before $U$ is chosen.
Its inventory is nonempty. The extra eligibility test is exactly
the conjunction of the original tests, not a relaxed test. The
new gas $\widehat G$ uses pairwise disjoint $Z_C$ and exact coverage
of $F(U)$ by the $J_C$.

\subsection{The exact compatible-family bijection}
The claimed identity is
\begin{equation}\label{eq:fusion-gas-identity}
                     \widehat G(U)=G(U)\exp\left(\sum_j\ell_j(U)\right).
\end{equation}
First expand the right side as a sum over a finite admissible hard
family $A$ and a finite set $B$ of logarithmic labels:
\begin{equation}\label{eq:fusion-expanded}
               \sum_{A\ {\rm admissible}}\sum_{B\ {\rm finite}}
                    \prod_{i\in A}K_i(U)\prod_{j\in B}(e^{\ell_j(U)}-1).
\end{equation}
To justify this, absolute summability of $\ell_j(U)$ implies absolute
summability of $e^{\ell_j(U)}-1$. Outside finitely many indices
$|\ell_j(U)|\le1$, and $|e^x-1|\le2|x|$ there. Finite products
$\prod_j[1+(e^{\ell_j}-1)]$ equal the corresponding exponentials;
their expansions converge absolutely, dominated by
$\prod_j(1+|e^{\ell_j}-1|)$. Passing to the limit proves the
logarithmic part of \eqref{eq:fusion-expanded}. The hard sum is
absolutely dominated by $\prod_i(1+|K_i|)$, so the double sum is
also absolutely convergent.

Consider a nonzero term of \eqref{eq:fusion-expanded} and split
$A\sqcup B$ into connected components of its support intersection
graph. Every component contains a hard label. In fact, if it contains
a logarithmic label $j$, its nonzero factor implies $\ell_j(U)\ne0$.
By \eqref{eq:fusion-vanishing}, some flagged site lies in $Y_j$.
Exact coverage of all flags by the old hard inventories places that
site in some $I_i$, $i\in A$. Since $I_i\subseteq X_i$, the labels
$i,j$ have overlapping supports and lie in the same component.
Each component is therefore an allowed new label.

Distinct components have disjoint supports. For a component $C$,
any flag in $Z_C$ belongs to an old inventory $I_i$; its hard support
then meets $Z_C$ and so $i$ lies in $C$. Conversely every $I_i$ in
$C$ consists of flags and lies in $Z_C$. Thus $Z_C\cap F(U)=J_C$.
Every old hard factor in the nonzero term is eligible. The product
of the new activities over the components consequently equals the
original term, and their inventories cover $F(U)$.

Conversely, take an admissible new family with nonzero product and
flatten its labels. Within each component the old hard supports
are disjoint by definition; between components they are disjoint
because the whole component supports are disjoint. The old hard
inventories cover $F(U)$, and the local eligibility indicators in
\eqref{eq:fusion-new-activity} equal one. Flattening therefore gives
an admissible hard family and a finite set of logarithmic labels.
No label is duplicated between components, because its support is
nonempty and would otherwise meet both disjoint component supports.
The connected components of the flattened support graph are exactly
the original new labels. This proves that flattening and taking
components are inverse on nonzero terms.

The identity follows by this weight-preserving bijection and absolute
convergence of \eqref{eq:fusion-expanded}. In particular the bijection
transports its absolute sum to that of the new gas. It also includes
the empty-family case: when there are no flags, every logarithm
vanishes and no nonzero new component can occur.

The new activities have the required locality and symmetry. Products
inherit gauge invariance; flags and all old eligibility predicates
are gauge invariant, so the two indicators are also invariant.
Measurability follows from measurable predicates and finite products.
For locality on $Z_C$, agreement of the configurations and their
flag data there gives agreement on every constituent support and
on every eligibility test. Finally the first indicator enforces
exact local inventory equality. These arguments apply to the
original addresses and predicates throughout.

\subsection{Preserving the derivative allowance}
Now let the hard factors be $K_i(\theta,U)$, with the logarithmic
terms independent of $\theta$. All addresses and eligibility tests
remain fixed. Let $D\ge0$, and suppose the hard factors have the
joint bound of \eqref{eq:leaf-joint-bound} with majorants $p_i$ on
$|\theta|D\le1$, and are independent of $\theta$ whenever their
support does not contain a fixed anchor $o$. Suppose also that
$|\ell_j(U)|\le v_j\le1$, and, for an exponent $u\ge0$,
\begin{equation}\label{eq:fusion-site-inputs}
 \sup_x\sum_{i:x\in X_i}p_i e^{(u+1)c(X_i)}\le b,
 \qquad
 \sup_x\sum_{j:x\in Y_j}v_j e^{(u+1)c(Y_j)}\le w.
\end{equation}
Choose $\kappa\ge0$ and define inflated constituent prices
\[
          a_i=e^\kappa p_i\quad(i\ {\rm hard}),\qquad
          a_j=2v_j\quad(j\ {\rm logarithmic}),\qquad
                              \Gamma=e^\kappa b+2w.
\]
For a logarithmic factor the constant joint bound is $2v_j$.
For a hard factor the joint bound is $e^{-\kappa}a_i$.
By the product rule for joint bounds, the product in a new label $C$
has joint bound
\[
 e^{-\kappa\#\{i\in C:\ i\ {\rm hard}\}}\prod_{v\in C}a_v
                           \le e^{-\kappa}\prod_{v\in C}a_v=:p_C.
\]
The two indicators in \eqref{eq:fusion-new-activity} are independent
of $\theta$, so the same bound holds for $\widehat K_C$.
It is independent of $\theta$ outside the anchor support as well:
if $o\notin Z_C$, each hard constituent is independent of $\theta$.

The constituent site price at exponent $u+1$ is at most $\Gamma$.
If $h_B\Gamma\le1$, the distinct connected-set estimate
\eqref{eq:ope-forest-price}, with reserve one, gives
\begin{equation}\label{eq:fusion-site-output}
       \sup_x\sum_{C:x\in Z_C}p_Ce^{u c(Z_C)}
                           \le e^{-\kappa}(e^\Gamma-1).
\end{equation}
Restricting to labels with nonempty hard part and separated hard
supports only decreases the nonnegative sum. Thus an allowance $A$
is sufficient whenever
\[
                    e^{-\kappa}(e^\Gamma-1)\le A.
\]
All the properties in \eqref{eq:leaf-season-bounds} are preserved,
with the same anchor and insertion debt $D$ and with new site price
$A$. Applying Section~\ref{sec:leaf-expansion} then gives marked
forced cost $DA$ on the fused alphabet. The original derivative
parameter is never rescaled when a logarithmic contribution is
absorbed.

For uniform absolute convergence in this argument, the finite lattice
and nonempty supports imply summability of $p_i$ and $v_j$ from
\eqref{eq:fusion-site-inputs}. The same estimate and
\eqref{eq:fusion-site-output} give summability of $p_C$. Thus the
scalar bijection, the parameter derivatives, and the site bounds
refer to the same series. Applying this construction successively
to the two localized logarithmic corrections proves the two gas
shift identities used in \eqref{eq:leaf-tilted-identity}, provided
their actual corrections satisfy \eqref{eq:fusion-vanishing}, the
fixed-address conditions and the purchased bounds above. Producing
those corrections from the nonlinear block integral, and paying these
inequalities on the selected history, remain the analytic inputs.
